\section{The Rendering Equation}\label{sec:rendering-equation}

The fundamental mathematical model for simulating light in physically-based rendering is the Rendering Equation,
introduced by Kajiya in 1986\cite{10.1145/15886.15902}.
It describes the equilibrium distribution of radiance in a scene.

\subsection{The Equation}\label{subsec:the-equation}

The Rendering Equation expresses the total spectral radiance $L_o$
leaving a point $x$ in direction $\omega_o$ as the sum
of emitted radiance $L_e$ and reflected radiance:

\begin{align}
    L_o(x, \omega_o) &= L_e(x, \omega_o) + \notag \\
    & \int_{\Omega} f_r(x, \omega_i, \omega_o) L_i(x, \omega_i) (\omega_i \cdot n) d\omega_i
    \label{eq:rendering-equation}
\end{align}

Where:
\begin{itemize}
    \item $L_o(x, \omega_o)$ is the total outgoing radiance.
    \item $L_e(x, \omega_o)$ is the emitted radiance (for light sources).
    \item $\Omega$ is the unit hemisphere centered at normal $n$ containing all incoming directions $\omega_i$.
    \item $f_r(x, \omega_i, \omega_o)$ is the \gls{brdf}, defining how light is scattered.
    \item $L_i(x, \omega_i)$ is the incoming radiance from direction $\omega_i$.
    \item $(\omega_i \cdot n)$ is the geometry term (Lambert's cosine law),
        accounting for the spreading of light over the surface area (often written as $\cos\theta_i$).
\end{itemize}

\paragraph{Usage of Radiance ($L$)}

Radiance ($L$) is the central quantity in rendering for several reasons:
\begin{enumerate}
    \item \textbf{Invariance}: Radiance is constant along a ray in a vacuum.
    This implies $L_i(x, \omega) = L_o(y, -\omega)$ where $y$ is the visible surface point.
    This property enables ray tracing: to find the radiance reaching the camera,
    we simply trace a ray into the scene and calculate the radiance leaving the intersected surface.
    \item \textbf{Sensor Response}: Camera sensors and the human eye respond to radiance
    (integrated over the pixel area and solid angle).
    \item \textbf{Light Transport}: It describes the flow of light energy through space.
\end{enumerate}

The recursive nature of $L_i$ (which is the $L_o$ from another point)
makes this integral equation difficult to solve analytically.
Therefore, most implementations use numerical methods like Monte Carlo integration (see \cref{sec:monte-carlo}).

\subsection{Spherical Integration}

To solve the integral numerically, we often parameterize the hemisphere $\Omega$
using spherical coordinates $(\theta, \phi)$.
The differential solid angle $d\omega_i$ changes variables as follows:
\begin{equation}
    d\omega = \sin\theta d\theta d\phi\label{eq:omega-cov}
\end{equation}

Substituting this into the rendering equation gives the explicit double integral form:

\begin{align}
    \int_{\Omega} (\dots) d\omega_i &= \int_{0}^{2\pi} \int_{0}^{\frac{\pi}{2}} L_i(x, \theta, \phi) \cdot \notag \\
    & f_r(x, \theta, \phi, \omega_o) \cos\theta \sin\theta d\theta d\phi
\end{align}

Here, $\theta$ is the polar angle (from the normal $n$, range $[0, \pi/2]$)
and $\phi$ is the azimuth angle (range $[0, 2\pi]$).
The term $\cos\theta$ comes from the geometry term $(\omega_i \cdot n)$,
and $\sin\theta$ comes from the differential solid angle measure.

\subsection{Color Spectrum}

Radiance $L$ varies with wavelength ($L(\lambda)$).
In computer graphics, strictly accurate rendering would require integrating over the entire visible spectrum.
However, for efficiency, we typically sample this spectrum at three discrete wavelengths corresponding to Red, Green,
and Blue (RGB).
This mimics the tristimulus response of the human eye.
While sufficient for most purposes, this approximation can fail for phenomena like
dispersion, fluorescence or metamerism resulting in inaccurate results.

\subsection{Energy Conservation (White Furnace Test)}

A physically-based material must conserve energy: it cannot reflect more light than it receives.
The ``White Furnace Test'' is a thought experiment (or unit-test) used to verify this.
If a non-absorbing object is placed in an environment with uniform radiance $L_i = 1$ everywhere,
the outgoing radiance $L_o$ must also be exactly $1$ (as long as no other form of light is emitted).
For a perfectly diffuse (Lambertian) surface with albedo $\rho$, the \gls{brdf} is constant $f_r = c$.
Integrating the cosine term over the hemisphere yields $\pi$:
\[
    \int_{\Omega} \cos\theta d\omega = \pi
\]
Therefore, to ensure $L_o = \rho L_i$, we must normalize the \gls{brdf}:
\[
f_{r, \text{diffuse}} = \frac{\rho}{\pi}
\]
The factor of $1/\pi$ is essential for energy conservation.